\section{Neutron Star Composition and Formation}
\label{sec:composition}
 
With the thermodynamic properties of the ideal Fermi gas established in
Section~\ref{sec:fermi}, it is now possible to examine the physical
conditions inside a neutron star and assess whether the free neutron
Fermi gas is a justified model for the stellar interior. This section
addresses three questions in turn: how neutron stars form, why their
interiors are composed predominantly of neutrons, and why those neutrons
constitute a highly degenerate Fermi gas to which the zero-temperature
approximation applies.
 
\subsection{Core-Collapse Supernovae and Neutron Star Formation}
 
Neutron stars are the remnants of massive stars with main
sequence masses $M \gtrsim 8\,M_\odot$ that end their lives in
core-collapse supernovae \cite{Heger2003}. Over the course of stellar
evolution, successive stages of nuclear burning build up an iron core,
which is the endpoint of exothermic fusion. When the iron core grows
beyond the effective Chandrasekhar mass of approximately $1.4\,M_\odot$,
electron degeneracy pressure can no longer support it and the core
collapses. Collapse halts when the core reaches nuclear saturation
density $\rho_0 \approx 2.3 \times 10^{14}$~g~cm$^{-3}$, at which
point the repulsive component of the nuclear force stiffens the equation
of state and the core bounces, launching the supernova shock that
disrupts the surrounding stellar envelope \cite{Burrows2021}. The
remnant is a proto-neutron star that cools rapidly through neutrino
emission, settling within seconds into a cold, degenerate object.

 
\subsection{Neutronization and the Composition of the Interior}
 
The composition of the neutron star interior is determined by the process
of neutronization. In ordinary stellar matter, neutrons are unstable against beta
decay $n \rightarrow p + e^- + \bar{\nu}_e$, with a free-neutron lifetime
of approximately 15 minutes. However, at the extreme densities of the
neutron star interior the reverse reaction,
%
\begin{equation}
    p + e^- \rightarrow n + \nu_e,
    \label{eq:neutronization}
\end{equation}
%
becomes energetically favorable. This occurs when the electron chemical
potential $\mu_e$ exceeds the neutron-proton mass difference,
%
\begin{equation}
    \mu_e > (m_n - m_p)c^2 = \Delta m\, c^2 = 1.293~\text{MeV}.
    \label{eq:neutronization-condition}
\end{equation}
%
From equation~\eqref{eq:fermi-momentum} of Section~\ref{sec:fermi}, the
electron Fermi momentum rises with electron number density as
$p_{F,e} \propto n_e^{1/3}$, and the corresponding ultra-relativistic
Fermi energy $E_{F,e} = p_{F,e}c$ therefore increases with density.
The condition \eqref{eq:neutronization-condition} is satisfied once the
density exceeds approximately $10^7$~g~cm$^{-3}$, well below the
densities found throughout the bulk of a neutron star
\cite{shapiro_teukolsky_1983}. The result is that electron capture drives the
composition toward a neutron-dominated state: in beta equilibrium, the
proton fraction satisfies $Y_p \lesssim 0.1$ throughout most of the
interior, and to a good first approximation the stellar matter may be
treated as a single-component gas of free neutrons. It is this
approximation that justifies the single-species Fermi gas model of
Section~\ref{sec:fermi}.
 
% ------------------------------------------------------------
\subsection{The Degeneracy Condition}
% ------------------------------------------------------------
 
The zero-temperature approximation adopted throughout Section~\ref{sec:fermi}
requires $T \ll T_F$, where the Fermi temperature $T_F = E_F/k_{\mathrm{B}}$
is set by the Fermi energy. Using equation~\eqref{eq:fermi-momentum},
the non-relativistic Fermi energy of the neutron gas is
%
\begin{equation}
    E_F = \frac{p_F^2}{2m_n}
        = \frac{\hbar^2}{2m_n}
          \left(3\pi^2 n\right)^{2/3},
    \label{eq:fermi-energy}
\end{equation}
%
where $n = \rho / m_n$ is the neutron number density. Evaluating at a
representative interior density $\rho \sim 2\rho_0 \approx
4.6 \times 10^{14}$~g~cm$^{-3}$ gives
%
\begin{equation}
    E_F \sim 60~\text{MeV},
    \label{eq:fermi-energy-numerical}
\end{equation}
%
corresponding to a Fermi temperature
%
\begin{equation}
    T_F = \frac{E_F}{k_{\mathrm{B}}} \sim 7 \times 10^{11}~\text{K}.
    \label{eq:fermi-temperature}
\end{equation}
%
The actual interior temperature of a neutron star older than a few
minutes is $T \sim 10^8$--$10^9$~K, set by neutrino cooling following
the proto-neutron star phase \cite{LattimerPrakash2004}. The ratio
%
\begin{equation}
    \frac{T}{T_F} \sim 10^{-3}\text{--}10^{-2}
    \label{eq:degeneracy-ratio}
\end{equation}
%
is small enough that thermal corrections to the equation of state are
negligible at the level of accuracy of the free Fermi gas model. The
zero-temperature approximation is therefore self-consistently justified,
and the step-function form of the Fermi-Dirac distribution used in
deriving equations~\eqref{eq:number-density}-\eqref{eq:pressure-rel}
is an excellent approximation.
 
% ------------------------------------------------------------
\subsection{The Relativistic Parameter and Its Implications}
% ------------------------------------------------------------
 
A further check concerns the degree of relativistic degeneracy. The
relevant dimensionless parameter is the ratio of the Fermi momentum to
the neutron rest momentum $m_n c$,
%
\begin{equation}
    x_F \equiv \frac{p_F}{m_n c}
             = \frac{\hbar}{m_n c}
               \left(3\pi^2 n\right)^{1/3}.
    \label{eq:relativistic-parameter}
\end{equation}
%
Evaluating at $\rho \sim \rho_0$ gives $x_F \sim 0.35$, rising to
$x_F \sim 0.7$ at $\rho \sim 2\rho_0$. Since $x_F$ is neither much
less than nor much greater than unity, the neutrons occupy the
\emph{intermediate relativistic regime}: neither the non-relativistic
limit $P \propto \rho^{5/3}$ nor the ultra-relativistic limit
$P \propto \rho^{4/3}$ is uniformly valid throughout the star. This
confirms that the full relativistic equation of state,
equations~\eqref{eq:energy-density-rel}--\eqref{eq:pressure-rel}, is
necessary for a quantitative treatment, as employed by Oppenheimer and
Volkoff \cite{Oppy}.
 
\subsection{Internal Structure and the Domain of the Model}
 
Neutron stars are not uniform objects: they possess a layered internal
structure whose composition varies with depth \cite{Chamel008}. The
outer crust consists of a lattice of neutron-rich nuclei in a
relativistic electron gas, the inner crust contains free neutrons
coexisting with nuclear clusters, and the inner core at densities above
$\sim 2\rho_0$ may harbor exotic phases whose equation of state is
poorly constrained. The free neutron Fermi gas is a reasonable
approximation only in the outer core, spanning roughly
$0.5\rho_0 \lesssim \rho \lesssim 2\rho_0$. This region constitutes the
bulk of the stellar volume and contains most of the mass, so the model
captures the dominant physics governing the global properties of the
star even though it is not valid everywhere.

